\section{Measurement Protocol}
\label{app:protocol}

This appendix specifies the protocol for the questions in \tabref{tab:results} in
enough detail to be executed by a party other than us. None of it has been executed
for this draft (\secref{sec:eval}). Where a choice would materially affect the
result, we state the choice and the reason, so that a different choice by a
replicator is visible as a deviation rather than absorbed as noise.

\subsection{Common setup}

\para{Provider configuration.} All measurements run with the deterministic fake
\code{LLMProvider}, \code{EvalProvider}, and \code{Promoter} unless a row explicitly
requires a real provider. This is the protocol's most important decision: it makes
E1 and E2 measure \sys rather than an inference service, and it makes results
reproducible across time, since a hosted model's latency is not a stable quantity.
Rows that need real providers (E3) say so.

\para{Topology.} Standalone (topology~B, \figref{fig:topologies}), so that \sys and
\mlflow{} are separate failure domains and \sys's own cost is separable. Report the
embedded topology separately if measured at all; do not average the two.

\para{Store configuration.} PostgreSQL~17 with pgvector and Apache~AGE; S3-compatible
object storage; \mlflow{} 3.14 or later with its own logical database. Record the
schema revision and the tree revision (\code{\statGitRev}) with every result.

\para{Warm-up and reporting.} Discard the first 60~seconds of every throughput or
latency run. Report p50, p95, and p99 --- never the mean alone, since the control
plane's tail is the operationally relevant quantity. Report the number of trials and
the interval, not a single figure.

\subsection{E1 --- Control-plane cost}

\para{E1(a), governed-read overhead.} Compare two clients issuing the same logical
operation: one resolving a \gasset through \sys, one reading a hard-coded version
from local configuration. Both must exercise the same downstream work, so the
difference isolates the control-plane path. Vary offered load until p99 degrades,
and report added latency at 50\% of the saturation point rather than at saturation ---
added latency measured at saturation is dominated by queueing and is not the number
an adopter needs.

\para{E1(b), alias resolution.} Measure resolution alone, with the downstream call
stubbed. Report cold (first resolution after process start) and warm separately;
reporting only the warm figure would hide whatever cache-miss cost exists.

\para{E1(c), throughput per replica.} Fix a p99 target, increase offered load until
the target is violated, and report the sustained rate below it. Report the target,
not just the rate: a throughput figure without a latency bound is uninterpretable.

\subsection{E2 --- Queue arbitration}

\para{E2(d), concurrent mutual exclusion under replication.} Enqueue a fixed
population of $J$ jobs whose handlers record their own invocation to a
uniquely-constrained table. Run $N \in \{1,2,4,8\}$ replicas against one database.
\textbf{Assert zero concurrent double-ownership}, not a low rate: this is a
correctness claim and a statistical result would not support it. Repeat with
induced contention by starting all replicas simultaneously against a full queue.

Note what this does \emph{not} test. Exclusion under contention is not
exactly-once execution across crashes, and the two diverge by loop
(\tabref{tab:loops}). Run E2(d) separately per loop and report its own semantics:
at-most-once for refinement and calibration, at-least-once for the workflow-run
worker, the knowledge-base worker, and webhook delivery. A single aggregate number
across loops with different recovery policies would be meaningless.

\para{E2(e), drain throughput versus replicas.} Same setup, measuring completion
rate. Expect sub-linear scaling and report where it flattens --- the flattening point
is the database arbitration limit and is the number that matters for capacity
planning.

\para{E2(f), recovery after a mid-claim kill.} Claim a job, kill the process
uncleanly ($\mathtt{SIGKILL}$, so no lease is released), and record what happens.
The expected outcome is loop-specific and the test must assert the right one:

\begin{itemize}
\item \textbf{WorkflowRun} --- the run is requeued on lease expiry and resumes from
      its last checkpoint. Measure the interval and check it against the configured
      lease timeout plus one sweep period.
\item \textbf{Refinement} --- the job is marked \code{failed} by the janitor and
      does \emph{not} resume. Assert that terminal state and measure time-to-reap.
      An earlier version of this protocol expected resumption here; that
      expectation was wrong about the implementation.
\item \textbf{Calibration} --- no lease exists, so the row remains \code{running}
      indefinitely. Assert that, and treat it as a defect to be fixed rather than a
      behaviour to be certified.
\end{itemize}

Separately, kill the process \emph{after} an external effect has been issued but
before its ledger entry settles, and assert the effect is surfaced as
\code{indeterminate} rather than silently retried or silently dropped.

\subsection{E3 --- The refinement path}

E3 requires real providers, and results are therefore provider-specific. Record the
model identifier and the date for every run.

\para{E3(g), signal-to-candidate latency.} Measure wall-clock from a verified
verification item to a candidate-ready job, broken down per stage
(\figref{fig:pipeline}). Report the breakdown, not only the total: a total dominated
by optimizer time and a total dominated by evidence assembly have entirely different
engineering implications.

\para{E3(h), gate agreement with expert judgement.} The sample size must come from
a power analysis rather than from a round number: fix the smallest $\kappa$ worth
detecting, the expected prevalence of withheld candidates, and the desired power,
and derive $n$ from those. Stratify by \family and by task, because a gate that
agrees well on prompts and poorly on workflows is not described by one pooled
figure. Seal the test set and select the release threshold \emph{before} it is
opened; a threshold chosen after looking is not a measurement.

Have at least three annotators who \emph{did not author the judges} label each
candidate independently, and adjudicate disagreements into an expert gold set
rather than resolving them by majority alone. Report Cohen's $\kappa$ against the
adjudicated label and Fleiss's $\kappa$ among annotators, both with confidence
intervals that account for clustering within family and
task~\cite{cohen1960kappa,fleiss1971kappa}. Report per-class sensitivity and
specificity separately: the cost of a false advance and the cost of a false
withhold are not symmetric, and a single agreement number hides which one the gate
is making. Report inter-annotator agreement even --- especially --- if it is low,
since a low value bounds what the gate could possibly agree with.

\para{E3(i), off-corpus generalization.} This is the measurement designed to be able
to embarrass the paper's own position (D4, \secref{sec:design}). Calibrate judges on
corpus $A$. Evaluate agreement on corpus $B$ drawn from a \emph{different} task
distribution --- not a held-out split of $A$, which would measure something else.
Report the agreement drop from $A$ to $B$. A small drop weakens the argument for a
mandatory human at \mode{Apply}; report it either way.

\subsection{E4 --- Release and recovery}

\para{E4(j), rollback latency.} Measure from the rollback request to the first
production call observably resolving the restored target. Include the client-side
resolution, since a rollback that has been recorded but not yet observed by clients
is not complete from the operator's point of view.

\para{E4(k), reconstructability.} The protocol's design matters more than its
mechanics. Perform $R \geq 20$ releases with varied paths --- some rolled back, some
superseded, some with failing gates that stopped a job. Then give an independent
party read access to \emph{every authoritative store}: the relational database,
object storage, \emph{and} \mlflow{} --- but nothing else: no access to the
original operator, no chat logs, no tickets. Ask them to state, per release, what
changed, what evidence justified it, who authorized it, and what the prior state
was. Score each of the four elements as recoverable or not, and report per-element
rather than as a single success rate --- the interesting failure is one element
being systematically unrecoverable. Any element unrecoverable for any release
falsifies the reconstructability claim in \secref{sec:overview}, and should be
reported as such rather than averaged away.

Including \mlflow{} in that access set is a correction rather than a convenience.
\mlflow{} owns prompt versions, aliases, traces, and assessments
(\figref{fig:dataflow}), so a protocol that withheld it could not reconstruct a
prompt release even in principle, and would have measured the access model rather
than the records. The stronger variant of this experiment --- and the one worth
running second --- withholds \mlflow{} and asks whether \sys can emit a
\emph{self-contained evidence bundle} that suffices on its own. That is a design
requirement this paper does not yet meet, and \secref{sec:future} records it as
work rather than as a property.

\subsection{E6 --- Baselines}

A control plane measured only against itself establishes nothing about whether it
was worth building. Every E1, E4, and E6 measurement is therefore run against four
baselines as well as \sys, on the same task set and by the same participants:

\begin{enumerate}
\item \textbf{Git plus CI.} Prompts in files, review by pull request, release by
      deployment. This is what most teams do and is the baseline \sys must beat to
      justify itself at all.
\item \textbf{\mlflow{} alone.} Prompt Registry versions and aliases, with
      \code{optimize\_prompts()} and evaluation run by
      hand~\cite{mlflowpromptregistry,mlflowoptimize}. This is the sharpest
      baseline, because it has most of \sys's mechanisms and none of its wiring;
      the delta it isolates is precisely the wiring.
\item \textbf{LangSmith}~\cite{langsmithprompts} and \textbf{Langfuse}
      \cite{langfuseprompts}, using commit tags and labels respectively.
\end{enumerate}

\para{E6(o), task time and correctness.} Give each participant the same three
tasks --- diagnose a reported failure, produce and release a fix, and roll it back
--- on each system, counterbalanced for order. Report time-to-complete and, more
importantly, the number of E4(k) elements a third party can later recover. A
baseline that is faster but leaves nothing recoverable is a different point on the
trade-off, not a worse system, and should be reported as such.

\para{E6(p), the \mlflow{}-alone delta.} Run E4(k) against baseline~2 in
isolation. The prediction \sys's design makes is that \emph{what changed}, \emph{what
the prior state was}, and \emph{who moved the alias} are recoverable from
\mlflow{} alone, and that \emph{what evidence justified it} is not, because the
evaluation result is not bound to the alias move. If all four elements turn out to
be recoverable, the paper's central value claim is substantially weakened and
should be reported as such.

\para{E6, applied to non-prompt families.} Baselines 2--4 cannot run the workflow,
tool, or MCP-server tasks at all, since those artifacts are not registry citizens
in any of them. Report that as a \emph{scope} result rather than as a score: the
comparison is not that \sys performs the task better but that the baseline has no
mechanism for it, which is a different and weaker kind of claim than a timing win.

\subsection{E7 --- Operators}

\secref{sec:arch:claim} argues that per-family guarantees are the honest factoring,
and \secref{sec:arch} names the resulting risk: an operator who sees the same
version-history panel in five places infers five equivalent rollback guarantees and
gets five different ones. That is an empirical claim about people and cannot be
settled by a table.

\para{E7(q), guarantee comprehension.} Recruit $\geq 12$ participants with
production LLM-operations experience and no prior exposure to \sys. After a fixed
onboarding, show each a version-history panel for each of five families and ask
what rolling back would do. Score against \tabref{tab:families}. Report per-family
accuracy, not a pooled figure: the interesting result is \emph{which} family is
misread, and the prediction is that tools --- which have no rollback at all --- are
misread most often. Report confidence intervals and the onboarding materials
verbatim, since comprehension is trivially inflated by better training.

A falsifying outcome here is not that \sys is slow but that its central design
position is unusable: if operators cannot state the guarantee, documenting it
per-family is not a mitigation, and the capability-interface design of
\secref{sec:arch:claim} becomes a requirement rather than future work.

\subsection{Fault injection}

E2(f) covers one fault. The following matrix is the minimum for a paper that claims
a recovery story, and each cell states the expected observable so a run either
matches it or does not:

\begin{itemize}
\item \textbf{Worker \code{SIGKILL}} --- per-loop outcome as in E2(f).
\item \textbf{Database failover mid-transaction} --- either no intent commits and
      no provider call begins, or a durable operation names the exact requested
      effect. Inject failure before and after each state commit.
\item \textbf{\mlflow{} unreachable during a release} --- the rotation fails and
      the operation remains \code{reconcile\_required}; reconciliation observes
      the prior target and settles it \code{failed}.
\item \textbf{\mlflow{} unreachable after a rotation but before the commit} ---
      the operation remains \code{applying} or \code{reconcile\_required};
      reconciliation observes the requested target and settles it \code{applied}.
\item \textbf{Object storage unavailable} --- evidence write fails closed; a
      release does not proceed on missing evidence.
\item \textbf{Clock skew across replicas} --- lease expiry does not permit two
      owners. Skew each replica by more than one lease period.
\item \textbf{Duplicate webhook delivery} --- the receiver's settlement path
      absorbs it; report the observed duplicate rate rather than asserting zero.
\end{itemize}

Report each cell as pass, fail, or not run. A blank cell is more informative than a
cell filled by argument.

\subsection{Deviations to report}

A replicator changing any of the following should say so, because each materially
affects comparability: the provider configuration (fake versus real); the topology;
the annotator pool's independence from judge authorship in E3(h); the distributional
distance between corpora $A$ and $B$ in E3(i); and whether E4(k)'s independent party
had any channel to the original operators; the baseline versions used in E6, since
these are live products; and the onboarding materials used in E7, since
comprehension results are not comparable across different training.
